
\section{有限型与凝聚层的更多美好性质}

我们现在来描述\trQCSheaf{}在具备某些有限性条件时的若干重要事实。

\textbf{14.3.A.（练习：$\scr{H}om$ 对有限表现的源表现合理；对照\mbox{练习~1.5.H}）}\quad (a) （「$\scr{H}om(\mathrm{fpr},\mathrm{qcoh})$ 是 $\mathrm{qcoh}$」）设 $\scr{F}$ 是 $X$ 上一个\trFinitelyPresented{}（\trQC）的层，而 $\scr{G}$ 是 $X$ 上一个\trQCSheaf。证明：$\scr{H}om(\scr{F},\scr{G})$ 是一个\trQCSheaf。提示：局部地取 $\scr{F}$ 的一个「\trFinitelyPresented」：在 $\mathscr{U}\subset X$ 上，$\OO^{\oplus m}\to\OO^{\oplus n}\to\scr{F}\to 0$。对它施加左正合函子（自 $\OO_X$-\trModule{}到 $\OO_X$-\trModule{}，\S2.6.4）$\scr{H}om_{\OO_X}(\cdot,\scr{G})$。（旁注：要找一个例子，使\trScheme{} $X$ 上的层 $\scr{F}$ 与 $\scr{G}$ 都\trQC{}而 $\scr{H}om(\scr{F},\scr{G})$ 不\trQC，可设 $A$ 是一个带单值化元 $t$ 的\trDVR，$X=\Spec A$，并设 $\scr{F}=\widetilde{M}$、$\scr{G}=\widetilde{N}$，其中 $M=\oplus_{i=1}^{\infty}A$、$N=A$。于是 $M_t=\oplus_{i=1}^{\infty}A_t$，而当然，$N_t=A_t$。考虑同态 $\phi\colon M_t\to N_t$，它把 $M_t$ 第 $i$ 个因子中的 $1$ 送到 $1/t^{i}$。那么 $\phi$ 不是 $\Hom_A(M,N)$ 中任何元素的局部化。）(b) （「$\scr{H}om(\mathrm{fpr},\mathrm{coh})$ 是 $\mathrm{coh}$」）若进而 $\scr{G}$ 是\trCoh{}的，证明 $\scr{H}om(\scr{F},\scr{G})$ 也\trCoh。(c)~设 $\scr{F}$ 是 $X$ 上一个\trFinitelyPresented{}的层，而 $\scr{G}$ 是 $X$ 上一个\trQCSheaf。证明：$\scr{H}om(\scr{F},\cdot)$ 给出一个左正合协变函子 $\mathit{QCoh}_X\to\mathit{QCoh}_X$ 与 $\mathit{Coh}_X\to\mathit{Coh}_X$，而 $\scr{H}om(\cdot,\scr{G})$ 给出一个从 $X$ 上\trFinitelyPresented{}的层到 $\mathit{QCoh}_X$ 的左正合反变函子（对照\mbox{练习~2.6.I}）。

\textbf{14.3.1.}\ \emph{\trCoherentSheaf{}的对偶。}\quad 由\mbox{练习~14.3.A(b)}，若 $\scr{F}$ 是\trCoh{}的，则其对偶 $\scr{F}^{\vee}:=\scr{H}om(\scr{F},\OO_X)$（见 \S2.3.4）也是。这推广了\mbox{练习~14.1.D}中\trVectorBundle{}之对偶的观念。你在那道练习里的论证经推广便可证明：总存在一个\trNatural{}的「双对偶」态射 $\scr{F}\to(\scr{F}^{\vee})^{\vee}$。与\trVectorBundle{}的情形不同，这并不总是一个\trIsomorphism。（例如，设 $\scr{F}$ 是 $\mathbb{A}^1=\Spec k[t]$ 上伴随于 $k[t]/(t)$ 的\trCoherentSheaf，证明 $\scr{F}^{\vee}=0$。）「双对偶」态射为\trIsomorphism{}的那些\trCoherentSheaf{}称为\textbf{\trReflexiveSheaf}（reflexive sheaves），但我们不会用到这一观念。典范映射 $\scr{F}\otimes\scr{F}^{\vee}\to\OO_X$ 称为\emph{\trTraceMap}（trace map）——你能看出这是为什么吗？

\textbf{14.3.B.（练习）}\quad 设
\begin{equation*}\tag{14.3.1.1}\label{eq:h-locally-free-sequence}
0\longrightarrow\scr{F}\longrightarrow\scr{G}\longrightarrow\scr{H}\longrightarrow 0
\end{equation*}
是概形 $X$ 上\trQCSheaf{}的一个正合列，其中 $\scr{H}$ 是一个\trLocallyFree{}的\trQCSheaf，并设 $\scr{E}$ 是一个\trQCSheaf。由 $\scr{H}om$ 的左正合性（\mbox{练习~2.6.I}），
\[
0\longrightarrow\scr{H}om(\scr{H},\scr{E})\longrightarrow\scr{H}om(\scr{G},\scr{E})\longrightarrow\scr{H}om(\scr{F},\scr{E})\longrightarrow 0
\]
是正合的，惟右端可能例外。证明它在右端也正合。（提示：这是局部的问题，故你可以假定 $X$ 是仿射的，比如说 $X=\Spec A$，且 $\scr{H}=\widetilde{A^{\oplus 1}}$，于是 \eqref{eq:h-locally-free-sequence} 可以写成 $0\to M\to N\to A^{\oplus 1}\to 0$。证明这个正合列分裂，于是我们可以把它写成 $N=M\oplus A^{\oplus 1}$，且写法与该正合列相容。）特别地，若 $\scr{F}$、$\scr{G}$、$\scr{H}$ 都\trCoh，且 $\scr{H}$ \trLocallyFree，则我们有凝聚层的一个正合列
\[
0\longrightarrow\scr{H}^{\vee}\longrightarrow\scr{G}^{\vee}\longrightarrow\scr{F}^{\vee}\longrightarrow 0.
\]

\textbf{14.3.C.（练习：「$\Supp(\operatorname{FT}\scr{F})$ 是闭的」）}\quad 设 $\scr{F}$ 是一个\trAbelianGroups{}层。(a)~证明：\trScheme{} $X$ 上一个\trFiniteType{}\trQCSheaf{}的\trSupport{}是一个\trClosedSubset。（提示：\mbox{练习~6.6.B(b)}。）(b)~证明：\trQCSheaf{}的\trSupport{}未必是闭的。（提示：若 $A=\mathbb{C}[t]$，则 $\mathbb{C}[t]/(t-a)$ 是一个支撑在 $a$ 处的 $A$-\trModule。考虑 $\oplus_{a\in\mathbb{C}}\mathbb{C}[t]/(t-a)$。当心：这个例子在把 $\oplus$ 换成 $\prod$ 时便不成立了。）

我们接下来来到\trNakayama{}引理的一个几何解读；正是由于它，\trNakayama{}引理才应当被视为一条几何事实（配上一个代数证明）。非正式地说：若一个\trFiniteType{}\trQCSheaf{}的某些\trSection{}「在 $\mathbf{p}$ 处」生成该层（即在 $\mathbf{p}$ 处的\trFiber{}上生成之），那么它们也在 $\mathbf{p}$ 「附近」生成（在 $\mathbf{p}$ 的一个货真价实的\trOpen\trNeighborhood{}中）。

\textbf{14.3.D.（有用练习：几何\trNakayama{}引理——对\trFiniteType{}的层而言，$\mathbf{p}$ 处的生成元便是 $\mathbf{p}$ 附近的生成元）}\quad 设 $X$ 是一个\trScheme，$\scr{F}$ 是一个\trFiniteType{}\trQCSheaf。证明：若 $\mathscr{U}\subset X$ 是 $\mathbf{p}\in X$ 的一个\trOpen\trNeighborhood，而 $a_1,\dots,a_n\in\scr{F}(\mathscr{U})$ 的像 $a_1|_{\mathbf{p}},\dots,a_n|_{\mathbf{p}}$ 生成\trFiber{} $\scr{F}|_{\mathbf{p}}$（其定义为 $\scr{F}_{\mathbf{p}}\otimes\kappa(\mathbf{p})$，\S4.3.8），则存在一个仿射开邻域 $\mathbf{p}\in\Spec A\subset\mathscr{U}$，使得「$a_1|_{\Spec A},\dots,a_n|_{\Spec A}$ 生成 $\scr{F}|_{\Spec A}$」，其确切含义如下：
\begin{enumerate}[label=(\roman*)]
\item $a_1|_{\Spec A},\dots,a_n|_{\Spec A}$ 作为 $A$-\trModule{}生成 $\scr{F}(\Spec A)$；
\item 对任意 $\mathbf{q}\in\Spec A$，$a_1,\dots,a_n$ 作为 $\OO_{X,\mathbf{q}}$-\trModule{}生成\trStalk{} $\scr{F}_{\mathbf{q}}$（从而对任意 $\mathbf{q}\in\Spec A$，诸\trFiber{} $a_1|_{\mathbf{q}},\dots,a_n|_{\mathbf{q}}$ 作为 $\kappa(\mathbf{q})$-\trVectorSpace{}生成\trFiber{} $\scr{F}|_{\mathbf{q}}$）。
\end{enumerate}
特别地，若 $\scr{F}|_{\mathbf{p}}=0$，则存在 $\mathbf{p}$ 的一个开邻域 $V$，使得 $\scr{F}|_V=0$。

\textbf{14.3.E.（有用练习：对\trFinitelyPresented{}的层而言，局部自由＝自由茎）}\quad 设 $\scr{F}$ 是\trScheme{} $X$ 上一个\trFinitelyPresented{}的层。证明：若 $\scr{F}_{\mathbf{p}}$ 对某个 $\mathbf{p}\in X$ 而言是一个自由的 $\OO_{X,\mathbf{p}}$-\trModule，则 $\scr{F}$ 在 $\mathbf{p}$ 的某个开邻域中\trLocallyFree。从而，$\scr{F}$ \trLocallyFree{}当且仅当 $\scr{F}_{\mathbf{p}}$ 对一切 $\mathbf{p}\in X$ 都是自由的 $\OO_{X,\mathbf{p}}$-\trModule。翻译过来便是：对\trFinitelyPresented{}的层而言，局部自由性是一个\trStalkLocal{}的性质，而自由茎蕴含附近的局部自由性。

扩充提示：找出 $\mathbf{p}$ 的一个开邻域 $\mathscr{U}$，以及生成 $\scr{F}_{\mathbf{p}}$ 的 $n$ 个属于 $\scr{F}(\mathscr{U})$ 的元素。利用几何\trNakayama{}引理（\mbox{练习~14.3.D}），证明这些\trSection{}对 $\mathscr{U}$ 内 $\mathbf{p}$ 的某个开邻域 $Y$ 中的一切 $\mathbf{q}$ 都生成 $\scr{F}_{\mathbf{q}}$。于是你便描述了一个满射 $\OO_Y^{\oplus n}\to\scr{F}|_Y$。证明这个映射的核是\trFiniteType{}的（利用引理~6.4.6，「有限表现者恒为有限表现」），从而具有闭支撑（比如说 $Z\subset Y$），并且该支撑不含 $\mathbf{p}$。于是 $\OO_{Y\setminus Z}^{\oplus n}\to\scr{F}|_{Y\setminus Z}$ 是一个\trIsomorphism。

这些说法从好几个角度来看都富有启发性。它表明：对\trFinitelyPresented{}的层而言，局部自由性是一个\trStalkLocal{}的条件。此外，在整概形上，任何\trFinitelyPresented{}的层 $\scr{F}$ 在\trGenericPoint{}处自动自由（你看出这是为什么了吗？），于是整概形上每个\trFinitelyPresented{}的层都在一个稠密开子集上\trLocallyFree。而且，在不可约概形的\trGenericPoint{}处为 $0$ 的\trFinitelyPresented{}层，必然在一个稠密开子集上为 $0$。最后两句话展示了一般点的效用；这类陈述在经典代数几何中本会显得更加神秘。一个表明有限性假设之必要性的例子：$\mathbb{Z}$-模 $A\subset\mathbb{Q}$，其中 $A$ 由分母\trSquareFree{}的分数组成；见 [Stacks, tag 065J]。

\textbf{14.3.F.（练习）}\quad （无挠层与挠层已在 \S6.1.4 中定义。）(a)~证明：正则（因而暗含局部\trNoetherian）\trCurve{}上无挠的\trCoherentSheaf{}都\trLocallyFree。(b)~证明：拟紧正则整\trCurve{}上的挠\trCoherentSheaf{}支撑在有限多个闭点处。(c)~设 $\scr{F}$ 是一条拟紧（为方便起见）正则\trCurve{}上的一个\trCoherentSheaf。描述一个典范的\trShortExactSequence
\[
0\longrightarrow\scr{F}_{\mathrm{tors}}\longrightarrow\scr{F}\longrightarrow\scr{F}_{\mathrm{lf}}\longrightarrow 0,
\]
其中 $\scr{F}_{\mathrm{tor}}$ 是一个挠层，而 $\scr{F}_{\mathrm{lf}}$ \trLocallyFree。

为解答\mbox{练习~14.3.F}，先用有用\mbox{练习~14.3.E}（局部自由性可在茎上检验）化归到\trDVR{}的情形，并回忆注记~6.4.4，即主理想整环 $A$ 上\trFinitelyGenerated{}\trModule{}的结构定理：任何这样的模都可以写成主模 $A/(a)$ 的直和。于是：

\textbf{14.3.2.\ \trhProposition{}。}\ ——\ \emph{若 $M$ 是\trPID{}上一个\trFinitelyGenerated{}的\trModule，则 $M$ \trTorsionFree{}当且仅当 $M$ 自由。}

（\mbox{练习~24.1.B} 与此密切相关。）

命题~14.3.2 在没有有限生成假设时不真：对一个适当一般的环 $A$，考虑 $M=K(A)$。若放弃主理想整环假设，它也不真：考虑 $(x,y)\subset\mathbb{C}[x,y]$ 与 $(x,y)\subset\mathbb{C}[x,y]/(xy)$。（这些例子需要一些验证。）这些例子表明：若放弃「维数 $1$」或「正则」假设，\mbox{练习~14.3.F(a)} 便不真。

\textbf{14.3.3.\ 一点处\trQCSheaf{}的\trFiber{}与\trRank。}

我们现在把 $\OO_X$-\trModule{}在一点 $\mathbf{p}\in X$ 处的\trFiber{}与\trRank{}的定义（\S4.3.8）应用到概形 $X$ 上的一个\trQCSheaf{} $\scr{F}$ 上。$\scr{F}$ 在 $\mathbf{p}$ 处的\emph{\trFiber}是\trVectorSpace
\[
\scr{F}|_{\mathbf{p}}:=\scr{F}_{\mathbf{p}}/\mathfrak{m}\scr{F}_{\mathbf{p}}=\scr{F}_{\mathbf{p}}\otimes_{\OO_{X,\mathbf{p}}}\kappa(\mathbf{p}),
\]
其中 $\mathfrak{m}$ 是 $\OO_{X,\mathbf{p}}$ 的\trMaximalIdeal，而 $\kappa(\mathbf{p})$ 一如通常是 $\mathbf{p}$ 处的\trResidueField{} $\OO_{X,\mathbf{p}}/\mathfrak{m}$。$\scr{F}$ 在一个含 $\mathbf{p}$ 的开集上的\trSection{}可以说在 $\mathbf{p}$ 处取值，其值是这个\trVectorSpace{}的一个元素。\trQCSheaf{} $\scr{F}$ 在一点 $\mathbf{p}$ 处的\emph{\trRank}是 $\dim_{\kappa(\mathbf{p})}\scr{F}_{\mathbf{p}}/\mathfrak{m}\scr{F}_{\mathbf{p}}$（可能无穷）。更明确地，在任何一处使 $\mathbf{p}=[\mathfrak{p}]$、$\scr{F}(\Spec A)=M$ 的仿射集 $\Spec A$ 上，秩为 $\dim_{\kappa(A_{\mathfrak{p}})}M_{\mathfrak{p}}/\mathfrak{p}M_{\mathfrak{p}}$。注意：这个\trRank{}的定义与\trLocallyFreeSheaf{}的\trRank{}观念是一致的。在局部自由的情形，\trRank{}是点的（局部）常值函数。其逆有时也真；见下文的\mbox{练习~14.3.K}。

\textbf{14.3.G.（练习）}\quad 考虑 $\mathbb{A}^1_k=\Spec k[t]$ 上对应于模 $k[t]/(t)$ 的\trCoherentSheaf{} $\scr{F}$。求出 $\scr{F}$ 在 $\mathbb{A}^1$ 的每一点处的\trRank。别忘了\trGenericPoint！

\textbf{14.3.H.（练习）}\quad 证明：在任意一点处，$\operatorname{rank}(\scr{F}\oplus\scr{G})=\operatorname{rank}(\scr{F})+\operatorname{rank}(\scr{G})$，并且 $\operatorname{rank}(\scr{F}\otimes\scr{G})=\operatorname{rank}\scr{F}\operatorname{rank}\scr{G}$。（提示：证明直和与张量积同环的商及局部化相交换，即：$(M\oplus N)\otimes_R(R/I)\cong M/IM\oplus N/IN$、$(M\otimes_R N)\otimes_R(R/I)\cong(M\otimes_R R/I)\otimes_{R/I}(N\otimes_R R/I)\cong M/IM\otimes_{R/I}N/IM$，等等。）

若 $X$ \trIrreducible，而 $\scr{F}$ 是 $X$ 上一个\trQC{}（通常是\trCoh）的层，则 $\operatorname{rank}\scr{F}$（不提及点）按约定指在\trGenericPoint{}处的秩。（例如，整概形上\trRank{}为 $0$ 的\trQCSheaf{}就是一个挠的\trQCSheaf；见定义~6.1.5。）

\textbf{14.3.I.（容易的练习）}\quad 证明：整概形 $X$ 上\trCoherentSheaf{}的\trRank{}在正合列中是可加的：若
\[
0\longrightarrow\scr{F}\longrightarrow\scr{G}\longrightarrow\scr{H}\longrightarrow 0
\]
是\trCoherentSheaf{}的一个正合列，证明
\[
\operatorname{rank}\scr{F}-\operatorname{rank}\scr{G}+\operatorname{rank}\scr{H}=0.
\]
提示：局部化是正合的。（当心：你的论证将用到「秩取在\trGenericPoint{}处」这一事实；$\mathbb{A}^1_k$ 上的例子
\[
0\longrightarrow\widetilde{k[t]}\overset{\times t}{\longrightarrow}\widetilde{k[t]}\longrightarrow\widetilde{k[t]/(t)}\longrightarrow 0
\]
表明闭点处的秩在正合列中不可加。）

若 $\scr{F}$ 是一个\trFiniteType{}\trQCSheaf，则 $\scr{F}$ 在 $\mathbf{p}$ 处的\trRank{}有限，而且由\trNakayama{}引理，这个\trRank{}就是把 $M_{\mathbf{p}}$ 作为 $A_{\mathbf{p}}$-\trModule{}看待时的最少生成元个数。

\textbf{14.3.J.（重要练习：对\trFiniteType{}\trQCSheaf{}而言，秩是\trUpperSemicontinuous{}的）}\quad 若 $\scr{F}$ 是 $X$ 上一个\trFiniteType{}\trQCSheaf，证明 $\operatorname{rank}(\scr{F})$ 是 $X$ 上的一个\trUpperSemicontinuous{}函数\footnote{原书此处作（a upper semicontinuous function），当系「an」之误，径改。——译者注}。提示：由几何\trNakayama{}引理（\mbox{练习~14.3.D}），在一点 $\mathbf{p}$ 处的生成元便是附近的生成元。（\mbox{练习~14.3.C} 中的例子表明\trFiniteType{}假设的必要性。）

\textbf{14.3.K.（重要练习：对既约概形上的\trFiniteType{}\trQCSheaf{}而言，常秩＝向量丛）}\quad (a)~设 $X$ \trReduced，$\scr{F}$ 是 $X$ 上一个\trFiniteType{}\trQCSheaf，且 $\scr{F}$ 的\trRank{}恒常，证明 $\scr{F}$ \trLocallyFree。

提示：在一个任意点 $\mathbf{p}$ 的某个开邻域中证明此结果，做法如下。设 $n=\operatorname{rank}_{\mathbf{p}}\scr{F}$。选取\trFiber{} $\scr{F}|_{\mathbf{p}}$ 的 $n$ 个生成元（即作为 $\kappa(\mathbf{p})$-\trVectorSpace{}的一组基）。由几何\trNakayama{}引理（\mbox{练习~14.3.D}），我们可以找到 $\mathbf{p}$ 的一个仿射开邻域 $\Spec A\subset X$，使得 $\scr{F}|_{\Spec A}=\widetilde{M}$，并且所选的 $\scr{F}|_{\mathbf{p}}$ 的那组生成元提升为 $M$ 的生成元 $m_1,\dots,m_n$。令 $\phi\colon A^{\oplus n}\twoheadrightarrow M$（$(r_1,\dots,r_n)\mapsto\sum r_i m_i$）为相应的满射。假设 $\ker\phi\neq 0$——比如说 $(r_1,\dots,r_n)\in\ker\phi$，其中 $r_i\neq 0$。由于 $X$ 既约，存在某个 $\mathfrak{p}$ 使 $r_i\notin\mathfrak{p}$。于是 $r_i$ 在 $A_{\mathbf{p}}$ 中可逆，故 $M_{\mathbf{p}}$ 的生成元少于 $n$ 个，与秩的恒常性矛盾。

(b)~然后利用秩的\trUpperSemicontinuous{}性（\mbox{练习~14.3.J}）证明：既约概形上的\trFiniteType{}\trQCSheaf{}在一个稠密开集上\trLocallyFree。(c)~证明：没有 $X$ 既约这一条件，(a) 可以不真。（提示：考虑 $A=k[x]/(x^2)$，$M=k$。）

\textbf{14.3.L.（练习）}\quad 设 $X$ \trReduced，而 $\phi\colon\scr{E}\to\scr{F}$ 是有限\trRank{}\trLocallyFreeSheaf{}之间的一个\trMorphism。证明：$\phi$ 是一个秩 $a$ 的\trVectorBundle{}映射（定义~14.2.4），当且仅当 $\phi|_{\mathbf{p}}\colon\scr{E}|_{\mathbf{p}}\to\scr{F}|_{\mathbf{p}}$ 对一切点 $\mathbf{p}\in X$ 都有秩 $a$。（提示：利用上一道练习与\mbox{练习~14.2.J}。）

你可以用\trRank{}这一观念来帮助直观地理解\trFiniteType{}\trQCSheaf，乃至\trQCSheaf。例如，\trFiniteType{}\trQCSheaf{}如下地推广了有限\trRank{}的\trVectorBundle：每一点处都有一个相伴的\trVectorSpace，尽管秩可以跳跃，它们却如人所愿地良好成族、彼此拼合。你不妨试着直观化例~14.3.G 中的那个例子。非既约性也能纳入这幅图像——你会如何描绘 $\Spec k[\epsilon]/(\epsilon^2)$ 上对应于 $k[\epsilon]/(\epsilon)$ 的那个\trCoherentSheaf（\mbox{练习~14.3.K(c)}）？那 $k[\epsilon]/(\epsilon^2)\oplus k[\epsilon]/(\epsilon)$ 又如何？

\textbf{14.3.4.}\ \emph{有限态射在一点处的\trDegree。}\quad 设 $\pi\colon X\to Y$ 是一个\trFinite\trMorphism。于是 $\pi_*\OO_X$ 是 $Y$ 上一个\trFiniteType{}（\trQC）的层，而这个层在一点 $\mathbf{p}$ 处的\trRank{}称为该\trFinite\trMorphism{}在 $\mathbf{p}$ 处的\textbf{\trDegree}（degree）。由\mbox{练习~14.3.J}，$\pi$ 的\trDegree{}是 $Y$ 上的一个\trUpperSemicontinuous{}函数。\trDegree{}可以跳跃：考虑把一个点嵌入一条直线中的那个闭嵌入，它对应于由 $t\mapsto 0$ 给出的 $k[t]\to k$。它也可以在你起初会觉得惊讶的情形中恒常——见\mbox{练习~10.3.4}，其中\trDegree{}总是 $2$，但这 $2$ 是以若干不同的方式取得的。

\textbf{14.3.M.（练习）}\quad 设 $\pi\colon X\to Y$ 是一个\trFinite\trMorphism。通过把定义解开，验证 $\pi$ 在 $\mathbf{p}$ 处的\trDegree{}是「$\mathbf{p}$ 的\trSchemeTheoreticPreimage{}上的函数所成空间」把\trResidueField{} $\kappa(\mathbf{p})$ 当作基域看待时的维数。特别地，\trDegree{}为零当且仅当 $\pi^{-1}(\mathbf{p})$ 为空。
